\documentclass[11pt,a4paper]{article}
\usepackage[margin=1in]{geometry}
\usepackage{hyperref}
\usepackage{enumitem}
\usepackage{graphicx}
\usepackage{xcolor}

% -----------------------------------------------------
% bvraghav-tiet-ucs503p-article.sty
% -----------------------------------------------------

% Variable: Lab Instructor
\makeatletter \newcommand{\BvrSetLabInstructor}[1]{%
  \gdef\Bvr@LabInstructor{#1}}
% default
\newcommand{\Bvr@LabInstructor}{Miss Nisha Thakur}
% public accessor
\newcommand{\BvrLabInstructorValue}{\Bvr@LabInstructor}
\makeatother

% Add subtitle
\usepackage{titling}
% -- Set title bold
\pretitle{\begin{center}\bfseries\fontsize{18bp}{18bp}\selectfont}
\makeatletter
\newcommand{\BvrSubtitle}[1]{%
  \posttitle{%
    \par\end{center}
    \begin{center}\large#1\end{center}
    \vskip 0.5em}%
}
\makeatother

% Hints in the section title(s)
\newcommand{\BvrHint}[1]{%
  {\normalfont\normalsize\itshape\hfill #1}}

% -----------------------------------------------------

\title{Git-Mind: A Self-Healing CI/CD Assistant with Human-in-the-Loop Review}

\BvrSetLabInstructor{Miss Nisha Thakur}

\BvrSubtitle{%
  Project Proposal for UCS503P  \\
  Submitted to: \BvrLabInstructorValue \\ \bfseries
  Thapar Institute of Engineering and Technology% 
}

\author{
  Mohit Srivastav\\ Roll No: \texttt{1024030368}
  \and
  Seerat\\ Roll No: \texttt{1024030361}
  \and
  Shreyas Tiwari\\ Roll No: \texttt{1024030384}
  \and
  Yuvraj Malik\\ Roll No: \texttt{1024030353}
}

\begin{document}
\maketitle

\section{Higher-order goal%
  \BvrHint{\ldots secondary framing}}
This project aims to reduce the engineering time lost to routine,
mechanically-fixable continuous-integration failures --- the syntax
slips, off-by-one errors, and stale logic that interrupt a
developer's flow without requiring deep judgement to resolve. While
developer productivity is a broad, organisation-wide concern, the
immediate engineering objective is to translate \emph{a narrow,
  well-verified class of automatic fixes} into a measurable,
deployable software system, rather than to claim general-purpose
autonomous debugging.

\section{Time-to-value %
\BvrHint{\ldots primary heuristic}}
\begin{itemize}[leftmargin=0.7cm]
    \item We validated the highest-risk assumption --- whether an
      LLM can reliably produce a working patch from an error log ---
      with a standalone script before building any surrounding
      infrastructure, and only proceeded once that assumption held
      up against a curated benchmark.
    \item We prioritize fast, falsifiable feedback over an
      impressive-looking demo: every fix generated by the agent is
      re-run against the project's own test suite before it is ever
      committed, so failure is caught before it is shipped rather
      than discovered later.
    \item Success in the first iteration is defined narrowly ---
      a working end-to-end loop from a detected bug to a real,
      human-reviewable pull request on a single-file class of bugs
      --- and is expanded only after that loop is demonstrated, not
      assumed.
\end{itemize}

\section{Problem Statement}
Developers routinely lose time to build and test failures that are
mechanically simple to fix but still require a full context switch
--- reading the log, locating the file, applying the correction, and
re-running the pipeline. Common pain points include:
\begin{itemize}[leftmargin=0.7cm]
    \item \textbf{Context-switch cost:} a one-line fix can cost
      several minutes of interrupted focus, disproportionate to the
      complexity of the bug itself.
    \item \textbf{Low visibility into automation attempts:} when
      teams do use auto-fix tooling, there is often no transparent
      record of what was tried, what passed verification, and what
      was discarded.
    \item \textbf{Unsafe automation:} naively deployed ``AI
      auto-fix'' tools risk committing plausible-looking but incorrect
      patches with no verification gate and no human checkpoint
      before code ships.
    \item \textbf{Repeated, avoidable toil:} the same narrow classes
      of bugs (missing imports, off-by-one errors, stale logic)
      recur across a codebase and are rarely worth a human's full
      attention individually.
\end{itemize}

\textbf{Why this matters now (contextual relevance):} As CI
pipelines run more frequently and teams ship smaller, faster
changes, the cumulative cost of routine failures grows even though
each individual failure is trivial to diagnose.

\section{Proposed Solution %
\BvrHint{\ldots Software Proposition}}
\subsection{Overview}
Build an \textbf{event-driven, human-in-the-loop} CI-failure
remediation system with the following components:
\begin{itemize}[leftmargin=0.7cm]
    \item \textbf{Failure ingestion:} a webhook receiver that
      captures CI failure events (repository, commit, error log)
      from GitHub.
    \item \textbf{Verification-gated fixer agent:} an LLM-driven
      worker that proposes a patch, applies it to a working copy,
      and re-runs the relevant test(s) --- a patch is committed
      \emph{only} if it passes.
    \item \textbf{Pull-request-based delivery:} verified fixes are
      pushed to a dedicated branch and opened as a normal pull
      request; the system never merges its own code.
    \item \textbf{Run logging:} every invocation --- successful or
      failed --- is recorded with its outcome, failure stage, and
      retry count, so the system's behaviour is auditable rather
      than anecdotal.
    \item \textbf{Live dashboard:} a node-graph visualization of
      repository activity and in-flight AI interventions.
\end{itemize}

\subsection{Core workflow}
\begin{enumerate}[leftmargin=0.7cm]
    \item A CI failure is detected (via webhook, or a manually
      triggered run during the current development phase) and
      queued with its error log and the implicated file.
    \item The fixer agent generates a candidate patch from the error
      context, retrying up to three times if generation or
      verification fails, feeding the new failure back into each
      subsequent attempt.
    \item The patch is applied to a fresh branch and the relevant
      test is executed locally; an unverified patch is discarded and
      the original file restored, with nothing committed.
    \item Only a verified patch is committed, pushed, and opened as
      a pull request for human review and merge.
    \item The dashboard and run log reflect the outcome regardless
      of whether the attempt succeeded.
\end{enumerate}

\subsection{Operational constraints for an educational setting}
\begin{itemize}[leftmargin=0.7cm]
    \item The fixer agent's scope is explicitly limited, on current
      evidence, to single-file, deterministic logic and syntax
      bugs; cross-file and stateful bug categories are tracked
      separately and not yet claimed as reliable.
    \item All git operations are bounded to a configured target
      repository path, with explicit guards against writing or
      pushing outside that scope.
    \item The system never auto-merges; a human reviewer is a
      mandatory step before any change reaches the main branch.
\end{itemize}

\section{Solution Approach %
\BvrHint{\ldots Engineering focus}}
This is an engineering course project: the contribution is a
verified, safety-gated pipeline built on existing LLM and developer
tooling, not a new machine-learning model or algorithm.

\subsection{Fix generation %
\BvrHint{\ldots non-research}}
Patch generation uses an off-the-shelf LLM API (Gemini, via
LangChain) prompted with the error log and file content, with output
parsed as a fenced code block rather than structured JSON --- a
deliberate choice made after empirically observing that JSON-mode
output was unreliable for multi-line code compared to a plain
code-fence extraction.

\subsection{Verification-first architecture %
\BvrHint{\ldots the core safety mechanism}}
Every candidate patch is subjected to the same local test the
original failure violated before it is allowed to reach git history.
This was validated directly: across an 8-case curated benchmark
spanning trivial, reasoning, cross-file, and stateful bug
categories, the verification gate correctly distinguished passing
patches from a patch that was semantically plausible but violated
the expected interface contract, without a human needing to inspect
the diff first.

\subsection{Event-driven backend}
\begin{itemize}[leftmargin=0.7cm]
    \item \textbf{Ingestion:} Express server handling GitHub
      webhooks with signature verification.
    \item \textbf{Queueing:} Redis-backed BullMQ workers, decoupling
      ingestion from the (slower, LLM-bound) fix-generation step.
    \item \textbf{Persistence:} MongoDB records of every run,
      shared between the API and worker services via a common
      schema module rather than direct filesystem coupling between
      services.
    \item \textbf{Live updates:} Socket.io pushes completed runs to
      a React Flow--based dashboard.
\end{itemize}

\section{Evaluation Criterion %
\BvrHint{\ldots measurable and attributable}}
Success is defined using metrics captured directly from the run log,
not from demo impressions.

\subsection{Primary evaluation metric}
\textbf{Verified-fix rate:} the percentage of triggered runs, within
the currently supported bug categories, that pass local verification
and result in an opened pull request.
\begin{itemize}[leftmargin=0.7cm]
    \item Current benchmark evidence: 7 of 8 cases in a curated,
      cross-category test set passed verification, with the single
      failure attributable to an unsupported bug category (stateful
      state-reset behaviour) rather than a verification-gate defect.
    \item Target for the pilot: $\ge$ 75\% verified-fix rate,
      measured across repeated runs (not a single pass) per bug
      category, reported separately by category rather than as one
      aggregate number.
    \item Attribution: measured directly from the persisted run log
      (outcome, attempt count, failure stage) for every invocation.
\end{itemize}

\subsection{Secondary evaluation metrics}
\begin{itemize}[leftmargin=0.7cm]
    \item \textbf{Time-to-PR:} elapsed time from failure detection
      to an opened pull request.
    \item \textbf{False-positive rate:} verified patches that pass
      the local test but are judged incorrect or unsafe on manual
      review (e.g., a patch that satisfies an assertion by
      narrowing behaviour rather than fixing it).
    \item \textbf{Category-wise reliability:} pass rate reported
      separately for trivial, single-file-reasoning, cross-file, and
      stateful bug categories, since these have shown materially
      different reliability in testing to date.
\end{itemize}

\subsection{Pilot validation plan %
\BvrHint{\ldots high-level}}
\begin{itemize}[leftmargin=0.7cm]
    \item Maintain and extend a curated benchmark repository of
      synthetic bugs across the categories above, each paired with
      an explicit assertion rather than a bare exception check.
    \item Run each benchmark case multiple times (not once) to
      report an actual pass rate rather than a single-run result,
      which testing so far has shown can be misleading.
    \item Compare manual fix time against agent-assisted time on a
      sample of real, injected CI failures once webhook-triggered
      automation is in place.
\end{itemize}

\section{Scalability %
\BvrHint{\ldots with foresight}}
\begin{enumerate}[leftmargin=0.7cm]
    \item \textbf{Scalable (theoretically):} The system is designed
      to support growth in concurrent repositories and failure
      volume by:
    \begin{itemize}[leftmargin=0.7cm]
        \item decoupling ingestion from processing via a queue,
          allowing worker instances to scale independently,
        \item sharing schemas between services through a common
          module rather than cross-service filesystem coupling, so
          services can be deployed and scaled independently,
        \item keeping the API layer stateless.
    \end{itemize}
    
    \item \textbf{Operationally deployable (practically):} The
      system is built entirely on containerizable, commonly hosted
      components:
    \begin{itemize}[leftmargin=0.7cm]
        \item a managed document database (MongoDB),
        \item a managed queue/cache (Redis),
        \item standard Node.js services deployable on common
          container hosting.
    \end{itemize}
\end{enumerate}

\section{Engine availability heuristic}
A mature set of ``engines'' (tooling/services) is available to
implement this system:
\begin{itemize}[leftmargin=0.7cm]
    \item Express for the webhook receiver and API layer.
    \item Redis and BullMQ for asynchronous job processing.
    \item MongoDB for run and repository state persistence.
    \item GitHub's REST API (via Octokit) and \texttt{simple-git}
      for repository operations.
    \item An LLM API (Gemini) via LangChain for patch generation.
    \item React, React Flow, and Socket.io for the live dashboard.
\end{itemize}
We use these mature components to focus engineering effort on the
verification gate, safety guards, and observability of the system,
rather than on infrastructure that already exists.

\section{Project scope and deliverables %
\BvrHint{\ldots iteration-friendly}}
\subsection{Initial deliverable %
\BvrHint{\ldots already demonstrated}}
\begin{itemize}[leftmargin=0.7cm]
    \item A dashboard frontend with a real, API-backed data layer
      (built first, per course guidance, ahead of the backend).
    \item A validated fixer-agent core loop: LLM-based patch
      generation, local test verification before commit, and
      automatic reversion of unverified patches.
    \item A real, manually triggered end-to-end run producing an
      actual GitHub pull request from a detected bug, with the
      diff manually confirmed to match the expected fix.
    \item Persistent logging of every run's outcome, retry count,
      and failure stage, decoupled from any single service's
      filesystem.
\end{itemize}

\subsection{Subsequent deliverables}
\begin{itemize}[leftmargin=0.7cm]
    \item GitHub webhook receiver with signature verification,
      replacing the current manually triggered flow.
    \item Redis/BullMQ-based asynchronous processing of incoming
      failure events.
    \item Live, socket-driven updates to the dashboard as fixes are
      attempted and resolved.
    \item A guard preventing AI-authored fix branches from
      re-triggering the fixer agent, to avoid self-referential
      retry loops.
    \item Expanded benchmark coverage and, where reliability
      permits, expanded scope into cross-file and stateful bug
      categories.
\end{itemize}

\section{Risks and mitigations}
\begin{itemize}[leftmargin=0.7cm]
    \item \textbf{LLM unreliability outside supported bug
      categories:} mitigated by the verification-before-commit gate
      (an unverified patch is never committed) and by explicitly
      scoping claims to categories with demonstrated pass rates
      rather than claiming general-purpose fixing.
    \item \textbf{Silent bad automation:} mitigated by mandatory
      human review before merge --- the system opens pull requests,
      it does not merge them.
    \item \textbf{Runaway retries / cost:} capped at three attempts
      per run, with non-retryable failures (missing configuration,
      safety-guard violations) short-circuiting immediately rather
      than exhausting all retries needlessly.
    \item \textbf{Unsafe file or repository access:} mitigated by
      explicit path-resolution checks that reject any write outside
      the configured target repository, and by scoping GitHub
      credentials to a single repository with the minimum required
      permissions.
    \item \textbf{Self-triggering fix loops:} planned mitigation is
      to exclude AI-authored branches from re-triggering webhook
      events, to be implemented alongside the webhook receiver.
    \item \textbf{Evaluation measurement ambiguity:} mitigated by
      logging outcome, attempt count, and failure stage for every
      run rather than relying on informal demo impressions.
\end{itemize}

\section{Summary}
Git-Mind proposes a verification-gated, human-in-the-loop system
that turns routine CI failures into reviewable pull requests rather
than autonomous merges. It meets the course criteria by validating
its highest-risk assumption before building surrounding
infrastructure, defining success through metrics captured directly
from a persisted run log rather than demo impressions, and scoping
its claims to the bug categories it has actually been shown to
handle reliably --- expanding that scope only as evidence supports
it.

\end{document}
